% GENERATED FILE. Edit canonical lessons, then run npm run generate:books.
% canonical-source-hash: fnv1a64:98ebe0224f37085e
% canonical-lessons: SA-C20-man, SA-C20-woman, SA-C20-child, SA-C20-old-person, SA-C20-people, SA-S121-letter-u

\chapter{People --- and the Four Noun Endings}
\label{ch:sa-people}
\hlchaptermodality{20}

\begin{chapteropening}
\textbf{By the end of this chapter:} \emph{I can name five kinds of person in Sanskrit and hear which of the four noun families each one belongs to.}

Complete SA-C20-people and demonstrate the chapter promise.
\end{chapteropening}

\section[naraḥ]{\sk{नरः} (naraḥ) — man, person}

\label{lesson:SA-C20-man}



\begin{quote}
Places, times, numbers — and nobody in them yet. This chapter supplies the people.
\end{quote}

\subsection*{You'll want to know: \sk{नरः}}
\textbf{\sk{नरः}} (\emph{naraḥ}) — \textquotedblleft{}man, person, human being.\textquotedblright{}

It is old enough to sit inside names you may already have heard: \emph{Narasiṃha}, \textquotedblleft{}man-lion,\textquotedblright{} and \emph{Nārāyaṇa}. Its Greek cousin is \emph{anēr} (genitive \emph{andros}), the element inside English \textbf{android} and the name \textbf{Andrew}.

Ending in \emph{-aḥ}, so you already know its family: masculine, like \sk{ग्रामः} and \sk{मार्गः}.

\begin{grammarlens}[title={What you've built}]
Your first person, and the ending told you what it was before the lesson did.
\end{grammarlens}

\subsection*{Your turn}
Say these aloud:

\begin{itemize}
\raggedright
  \item \emph{naraḥ} — \textquotedblleft{}man\textquotedblright{}
  \item \textquotedblleft{}a good man\textquotedblright{} — \emph{uttama naraḥ}
  \item the English word built on its Greek cousin (\emph{android})
\end{itemize}

\begin{tcolorbox}[breakable,colback=teal!4,colframe=teal!35!black,title={Before you move on}]
Which family does \sk{नरः} belong to, and how do you know? (Masculine — the \emph{-aḥ}.) Which English words carry its Greek cousin? (\emph{Android}, \emph{Andrew}.)
\end{tcolorbox}
\section[nārī]{\sk{नारी} (nārī) — woman}

\label{lesson:SA-C20-woman}



\begin{quote}
And the ending you have not met yet.
\end{quote}

\subsection*{You'll want to know: \sk{नारी}}
\textbf{\sk{नारी}} (\emph{nārī}) — \textquotedblleft{}woman.\textquotedblright{}

Look at it beside \sk{नरः} and the relationship is visible: the same \emph{nar-}, a different ending. Sanskrit builds a great many feminine nouns with a long \textbf{-\sk{ई}} (\emph{-ī}), and this is your first.

So the map has grown again: \emph{-am} neuter, \emph{-aḥ} masculine, \emph{-iḥ} a third family, and now \emph{-ī} feminine. Four endings, four families. You are not being asked to produce them — only to hear them, because everything Sanskrit does with a noun starts here.

\begin{grammarlens}[title={What you've built}]
A second person and a fourth ending. The system is bigger than two, exactly as the night lesson warned.
\end{grammarlens}

\subsection*{Your turn}
Say these aloud:

\begin{itemize}
\raggedright
  \item \emph{nārī} — \textquotedblleft{}woman\textquotedblright{}
  \item it against \emph{naraḥ}, and name the shared piece
  \item the four endings you have now met
\end{itemize}

\begin{tcolorbox}[breakable,colback=teal!4,colframe=teal!35!black,title={Before you move on}]
What do \sk{नरः} and \sk{नारी} share? (The stem \emph{nar-}.) Which ending marks this feminine? (Long \emph{-ī}.)
\end{tcolorbox}
\section[bālaḥ]{\sk{बालः} (bālaḥ) — child, boy}

\label{lesson:SA-C20-child}



\begin{quote}
The youngest person in the chapter.
\end{quote}

\subsection*{You'll want to know: \sk{बालः}}
\textbf{\sk{बालः}} (\emph{bālaḥ}) — \textquotedblleft{}child, boy, a young one.\textquotedblright{}

It carries a sense wider than English \textquotedblleft{}boy\textquotedblright{}: \emph{bāla} is anything in its early state, and the word is used of young animals and of the newly-risen sun as readily as of a child. The feminine \textbf{\sk{बाला}} (\emph{bālā}) is a girl.

Notice that the feminine here is \emph{-ā}, not the \emph{-ī} you just met. Sanskrit has more than one way to make a feminine, which is inconvenient and true.

\begin{grammarlens}[title={What you've built}]
Three people, and an honest complication: the feminine endings are not a single rule.
\end{grammarlens}

\subsection*{Your turn}
Say these aloud:

\begin{itemize}
\raggedright
  \item \emph{bālaḥ} — \textquotedblleft{}child\textquotedblright{}
  \item the feminine — \emph{bālā}
  \item the two different feminine endings you have now seen
\end{itemize}

\begin{tcolorbox}[breakable,colback=teal!4,colframe=teal!35!black,title={Before you move on}]
What does \sk{बाल} cover beyond \textquotedblleft{}boy\textquotedblright{}? (Anything in its early state.) Which two feminine endings have you met? (\emph{-ī} and \emph{-ā}.)
\end{tcolorbox}
\section[vṛddhaḥ]{\sk{वृद्धः} (vṛddhaḥ) — an old person, an elder}

\label{lesson:SA-C20-old-person}



\begin{quote}
The oldest person in the chapter — and a word that is careful in a way English is not.
\end{quote}

\subsection*{You'll want to know: \sk{वृद्धः}}
\textbf{\sk{वृद्धः}} (\emph{vṛddhaḥ}) — \textquotedblleft{}an old person, an elder, one who has grown.\textquotedblright{}

You already have \textbf{\sk{पुरातन}} (\emph{purātana}), which was glossed \textquotedblleft{}old — of a thing, not of a person.\textquotedblright{} Here is the other half of that pair, and the distinction is the point.

\emph{\sk{पुरातन}} is old as in \emph{ancient}: it has existed a long time. \emph{\sk{वृद्धः}} is old as in \emph{grown}: from the root \textbf{√\sk{वृध्}} (\emph{vṛdh}), \textquotedblleft{}to grow, increase.\textquotedblright{} An elder is not a thing that has lasted; an elder is someone who has grown.

English flattens both into \textquotedblleft{}old\textquotedblright{} and loses the compliment.

\begin{grammarlens}[title={What you've built}]
Four people — and a pair of words that shows why a one-word gloss is sometimes a lie by omission.
\end{grammarlens}

\subsection*{Your turn}
Say these aloud:

\begin{itemize}
\raggedright
  \item \emph{vṛddhaḥ} — \textquotedblleft{}an elder\textquotedblright{}
  \item which of \emph{purātana} and \emph{vṛddha} you would use for a temple
  \item what √\sk{वृध्} means (to grow)
\end{itemize}

\begin{tcolorbox}[breakable,colback=teal!4,colframe=teal!35!black,title={Before you move on}]
What root is \sk{वृद्धः} built on? (√\sk{वृध्}, \textquotedblleft{}to grow.\textquotedblright{}) Which word would you use of an old building? (\emph{Purātana}.)
\end{tcolorbox}
\section[janaḥ]{\sk{जनः} (janaḥ) — people, folk}

\label{lesson:SA-C20-people}



\begin{quote}
All of them at once.
\end{quote}

\subsection*{You'll want to know: \sk{जनः}}
\textbf{\sk{जनः}} (\emph{janaḥ}) — \textquotedblleft{}people, folk, a person, a race.\textquotedblright{}

From the root \textbf{√\sk{जन्}} (\emph{jan}), \textquotedblleft{}to be born, to beget\textquotedblright{} — the same root behind \emph{janma}, \textquotedblleft{}birth.\textquotedblright{} Its Latin cousin is \emph{gens} (\textquotedblleft{}clan, people\textquotedblright{}) and \emph{genus}, and through them English has \textbf{genus}, \textbf{gene}, \textbf{generate}, \textbf{genesis} and \textbf{nation} (from \emph{natio}, \textquotedblleft{}a being born\textquotedblright{}).

So \sk{जनः} does not mean a crowd. It means \emph{those who were born} — kinship rather than quantity.

\begin{grammarlens}[title={What you've built}]
Five people: man, woman, child, elder, folk — and every one of them arrived with its ending legible and a European cousin attached.
\end{grammarlens}

\subsection*{Your turn}
Say these aloud:

\begin{itemize}
\raggedright
  \item \emph{janaḥ} — \textquotedblleft{}people\textquotedblright{}
  \item three English words from its Latin cousin
  \item all five people, in order
\end{itemize}

\begin{tcolorbox}[breakable,colback=teal!4,colframe=teal!35!black,title={Before you move on}]
What does √\sk{जन्} mean? (To be born.) Which English words share it? (\emph{Gene}, \emph{genesis}, \emph{generate}, \emph{nation}.)
\end{tcolorbox}
\section[u]{\sk{उ} — one character, met inside words you already say}

\label{lesson:SA-S121-letter-u}



\begin{quote}
Before the new one: \sk{ज} — what does it do?

One character this time. One only — and you have been saying it for pages without knowing which mark on the page it was.
\end{quote}

\subsection*{Script you'll notice: \sk{उ}}
\textbf{\sk{उ}} — \emph{u}.

It is an \textbf{independent vowel} — the shape the vowel \emph{u} takes when a word begins with it, rather than the sign it becomes inside a word.

What it is made of:

\begin{itemize}
\raggedright
  \item a joined upper bowl and lower loop
  \item the top shirorekhā
\end{itemize}

You already say these, and every one of them has \sk{उ} somewhere inside it:

\begin{itemize}
\raggedright
  \item \textbf{\sk{उत्तम}} \emph{uttama} — good, best
  \item \textbf{\sk{उदक}} \emph{udaka} — water — the other word for it
  \item \textbf{\sk{उपसर्ग}} \emph{upasarga} — a prefix, the piece that turns \sk{गच्छति} into \sk{आगच्छति}
\end{itemize}

You have read the \emph{u} sound many times without ever meeting this shape, because inside a word it is a small hook tucked under the consonant instead. \sk{उ} is what it looks like when a word opens on it and there is no consonant to sit under.

\subsection*{Writing: \sk{उ}}
\begin{itemize}
\raggedright
  \item \textbf{1.} start at the headline junction, curve down and left around the upper bowl, then sweep back through the waist and around the lower loop without lifting
  \item \textbf{2.} lift and draw the top shirorekhā left-to-right
\end{itemize}

\textbf{Pen lifts: 1.} The pen comes up 1 times and no more.

\begin{quote}
Verified modern printed teaching form; everyday handwriting may simplify the curves.
\end{quote}

\begin{quote}
This is one attested teaching order and not a national standard — handwriting here is taught with school-to-school variation. Source: Saurmandal, ‘Devanagari \sk{उ} stroke order.svg’, panels 1–2, Wikimedia Commons, 5 August 2023.
\end{quote}

\subsection*{Your turn}
\begin{itemize}
\raggedright
  \item \emph{Look:} at these words, and find \sk{उ} in the ones that have it
\end{itemize}

\begin{quote}
\sk{उत्तम}  ·  \sk{नरः}
\end{quote}

\begin{itemize}
\raggedright
  \item \emph{Trace it:} \sk{उ} three times, saying \emph{u} as you finish each one
  \item \emph{Look:} back at any page of this chapter and find \sk{उ} once more
\end{itemize}

\begin{tcolorbox}[breakable,colback=teal!4,colframe=teal!35!black,title={Before you move on}]
Which character is this — \sk{उ}? What sound does it carry? (\emph{\textbf{u}}.) Name one word you already say that contains it.
\end{tcolorbox}
